\section{Limitations}
\textbf{Output-level evidence.} We analyse generated text. Stated assumptions need not reflect the computation that produced a decision, and our conclusions concern observable behaviour under one protocol.
\textbf{Synthetic cases and simple policies.} The domains are controlled probes. Real institutions use richer data and procedures, and our policies are illustrative thresholds, not professional standards.
\textbf{Lexical detectors.} The v2 rules are transparent but lexical; their error rates are reported (Sec.~\ref{sec:validation}) and every rate inherits them.
\textbf{Models.} Five open-weight models of 7--14B parameters, quantised as distributed by Ollama. Larger or proprietary models may behave differently.
\textbf{Prompting.} Requiring an explicit assumptions list may itself elicit assumptions; we measure sensitivity to prompt wording but not to removing the field.
\textbf{Disclosure.} Four of five models list assumptions in a minority of runs, so their assumption-level rates rest on few runs, and the annotated sample is dominated by Mistral and DeepSeek-R1 outputs, which make up most of the stated assumptions.
\textbf{Confidence.} Verbalized confidence is not a calibrated probability; we use it relationally.
